\documentclass[UTF8,11pt,a4paper,fontset=none]{ctexart}
\usepackage[left=21mm,right=21mm,top=22mm,bottom=23mm,headheight=15pt]{geometry}
\usepackage{fontspec}
\setmainfont{Times New Roman}
\setsansfont{Arial}
\setmonofont{Menlo}[Scale=0.90]
\setCJKmainfont{Songti SC}
\setCJKsansfont{Heiti SC}
\setCJKmonofont{Heiti SC}
\usepackage{amsmath,amssymb}
\usepackage{xcolor,graphicx}
\usepackage{booktabs,array,tabularx,longtable}
\usepackage{enumitem}
\usepackage{listings}
\usepackage{tikz}
\usetikzlibrary{arrows.meta,positioning,calc}
\usepackage{fancyhdr}
\usepackage{needspace}
\usepackage{hyperref}
\definecolor{Ink}{HTML}{172C44}
\definecolor{Teal}{HTML}{087F8C}
\definecolor{Muted}{HTML}{5B6C7B}
\definecolor{Pale}{HTML}{F2F6F9}
\definecolor{Rule}{HTML}{CCD9E3}
\definecolor{Rust}{HTML}{B45831}
\definecolor{CodeGreen}{HTML}{376D48}
\hypersetup{
  colorlinks=true,linkcolor=Teal,urlcolor=Teal,citecolor=Teal,
  pdftitle={LLVM IR 语法与实践入门：面向 Rx 编译器},
  pdfauthor={Rx Compiler 学习资料},
  pdfsubject={LLVM 22、RV32IM、LLVM IR 语法、项目实践}
}
\urlstyle{same}
\ctexset{
  section={format=\Large\sffamily\bfseries\color{Ink},beforeskip=18pt,afterskip=10pt},
  subsection={format=\large\sffamily\bfseries\color{Ink},beforeskip=13pt,afterskip=7pt},
  subsubsection={format=\normalsize\sffamily\bfseries\color{Teal},beforeskip=10pt,afterskip=5pt}
}
\setlength{\parindent}{0pt}
\setlength{\parskip}{4pt}
\linespread{1.13}
\clubpenalty=10000
\widowpenalty=10000
\setlist[itemize]{leftmargin=1.6em,itemsep=3pt,topsep=4pt}
\setlist[enumerate]{leftmargin=1.8em,itemsep=3pt,topsep=4pt}
\setlength{\emergencystretch}{2em}
\renewcommand{\arraystretch}{1.20}
\newcolumntype{Y}{>{\raggedright\arraybackslash}X}
\newcommand{\code}[1]{\texttt{#1}}
\newcommand{\lesson}[2]{\clearpage\section{#1}\label{#2}}
\newcommand{\aim}[1]{\noindent{\color{Teal}\sffamily\bfseries 本课目标}\quad #1\par\medskip}
\newcommand{\takeaway}[1]{\par\medskip\noindent\colorbox{Pale}{\parbox{\dimexpr\linewidth-2\fboxsep\relax}{\small\textbf{\color{Teal}记住这一点}\quad #1}}\par\medskip}
\newcommand{\projectnote}[1]{\par\smallskip{\small\textbf{\color{Rust}项目约定}\quad #1}\par\smallskip}
\newcommand{\checkq}[1]{\par\smallskip{\small\textbf{\color{Teal}自测}\quad #1}\par\smallskip}
\newcommand{\srcref}[2]{\par{\footnotesize\color{Muted}进一步核对：\href{#1}{#2}。}\par}
\pagestyle{fancy}
\fancyhf{}
\fancyhead[L]{\small\sffamily\color{Muted}Rx Compiler / LLVM IR}
\fancyhead[R]{\small\sffamily\color{Muted}语法与实践入门}
\fancyfoot[L]{\footnotesize\sffamily\color{Muted}LLVM 22 \quad RV32IM / ILP32}
\fancyfoot[R]{\small\sffamily\color{Ink}\thepage}
\renewcommand{\headrulewidth}{0.35pt}
\renewcommand{\footrulewidth}{0pt}
\lstdefinelanguage{LLVM}{
 morekeywords={target,datalayout,triple,source_filename,private,internal,constant,global,declare,define,
 ret,call,alloca,load,store,getelementptr,add,sub,mul,sdiv,udiv,srem,urem,
 and,or,xor,shl,ashr,lshr,icmp,eq,ne,slt,sle,sgt,sge,ult,ule,ugt,uge,
 zext,sext,trunc,br,label,phi,select,unreachable,align,to,inbounds,nsw,nuw,exact},
 morekeywords=[2]{void,i1,i8,i32,i64,ptr,true,false,null},
 sensitive=true,morecomment=[l]{;},morestring=[b]"
}
\lstdefinelanguage{Rx}{
 morekeywords={fn,let,mut,if,else,while,loop,break,continue,return,struct,impl,const,as,true,false},
 morekeywords=[2]{i32,u32,isize,usize,bool,Box,Vec},
 sensitive=true,morecomment=[l]{//},morestring=[b]"
}
\lstset{
 language=LLVM,
 basicstyle=\ttfamily\fontsize{9}{11.4}\selectfont,
 keywordstyle=\bfseries\color{Ink},
 keywordstyle=[2]\color{Teal},
 commentstyle=\itshape\color{Muted},
 stringstyle=\color{CodeGreen},
 numbers=left,numberstyle=\ttfamily\fontsize{7}{9}\selectfont\color{Muted},
 numbersep=8pt,xleftmargin=16pt,framexleftmargin=4pt,
 frame=tb,framerule=0.35pt,rulecolor=\color{Rule},
 backgroundcolor=\color{Pale},
 aboveskip=8pt,belowskip=8pt,
 breaklines=true,breakatwhitespace=false,
 columns=fullflexible,keepspaces=true,
 showstringspaces=false,tabsize=2,
 captionpos=b
}
\AddToHook{env/lstlisting/before}{\par\noindent\begin{minipage}{\linewidth}}
\AddToHook{env/lstlisting/after}{\end{minipage}\par}
\let\originalinputlisting\lstinputlisting
\renewcommand{\lstinputlisting}[2][]{%
  \par\noindent\begin{minipage}{\linewidth}%
  \originalinputlisting[#1]{#2}%
  \end{minipage}\par}
\tikzset{
 flowbox/.style={draw=Rule,fill=Pale,rounded corners=2pt,minimum height=9mm,
   text width=28mm,align=center,font=\small\sffamily},
 flowarrow/.style={-{Latex[length=2mm]},thick,color=Teal}
}
\setcounter{tocdepth}{1}

\begin{document}
\hypersetup{pageanchor=false}
\begin{titlepage}
\thispagestyle{empty}
{\sffamily\color{Teal}\large RX COMPILER \hfill STUDY NOTES / 01}
\vspace{26mm}

{\sffamily\bfseries\color{Ink}\fontsize{34}{43}\selectfont LLVM IR\par}
\vspace{3mm}
{\sffamily\bfseries\color{Ink}\fontsize{25}{34}\selectfont 语法与实践入门\par}
\vspace{9mm}
{\Large 面向 Rx 编译器的中文学习讲义\par}
\vspace{5mm}
{\color{Muted}从读懂一条指令，到写出可运行的模块。\\
以 LLVM 22 的 \code{ptr} 语法为基准，目标为 RV32IM / ILP32。}
\vspace{17mm}
\begin{lstlisting}[numbers=none,basicstyle=\ttfamily\fontsize{11}{15}\selectfont]
define i32 @add(i32 %a, i32 %b) {
entry:
  %sum = add i32 %a, %b
  ret i32 %sum
}
\end{lstlisting}
\vspace{9mm}
\begin{tabularx}{\linewidth}{@{}>{\bfseries\color{Teal}}l Y@{}}
适合谁 & 已完成 AST 与 semantic，准备实现 IR 生成的同学。\\
学什么 & 类型、值、地址、指令、函数、基本块、控制流和聚合寻址。\\
怎么学 & 每课先读示例，再改常量或控制流，最后运行并解释结果。\\
配套材料 & 12 个完整 \code{.ll} 示例、预期输出、验证脚本与 LaTeX 源文件。\\
\end{tabularx}
\vfill
{\color{Rule}\rule{\linewidth}{0.8pt}}
{\small\sffamily\color{Muted}项目学习版 \hfill 2026 年 11 月 2 日}
\end{titlepage}
\pagenumbering{arabic}
\hypersetup{pageanchor=true}
\section*{如何使用这份讲义}
\addcontentsline{toc}{section}{如何使用这份讲义}
这份资料解释“LLVM IR 代码如何阅读和书写”。工程中的模块划分、Builder 接口和 verifier 设计，请结合项目的 \code{docs/ir.md} 阅读。这里的每个例子都是为学习重新编写的短程序，语言规则与本项目的表示选择会明确区分。

\begin{itemize}
 \item \textbf{第一遍：}第 1 至 4 课，完成打印、函数调用、变量读写。
 \item \textbf{第二遍：}第 5 至 7 课，掌握运算、分支、短路和循环。
 \item \textbf{第三遍：}第 8 至 10 课，处理数组、结构体、引用及项目中的聚合传递。
 \item \textbf{动手验收：}第 11 课提供工具链；第 12 课给出练习和答案。附录可当指令速查表。
\end{itemize}

\projectnote{本项目第一版用栈槽保存源变量及分支结果。第 7 课的 \code{phi} 用于理解 SSA，是可选拓展；当前核心指令枚举尚未包含它。示例能够由 LLVM 接受，不等于当前学生编译器已实现对应特性。}

\textbf{代码怎么找。}完整例子位于本目录的 \code{examples/}。讲义中只展示部分函数时会说明；语法模板和错误示例不能直接当成完整模块运行。每个完整示例都有独立的入口和打印支持，不需要借用项目参考 Rust 编译器。

\textbf{验证状态。}配套脚本记录编译结果与精确输出，并明确区分本机模式和目标模式。本机运行时会在生成的副本中去掉 RV32 目标行；只有 Clang 22 交叉编译并在 REIMU 运行，才能称为项目目标链路通过。具体本次结果见同目录的 \code{build/validation-report.json}。

\clearpage
\tableofcontents
\lesson{先读懂一个完整模块}{sec:first}
\aim{认识模块、全局数据、函数、基本块和指令，能解释打印 42 的完整执行路径。}

LLVM IR 是带类型的中间表示。文本形式通常保存为 \code{.ll}。它已经把源语言中的“表达式和语句”拆成了更明确的计算、读写与跳转，但还没有指定某个值必须放进哪一个 RISC-V 物理寄存器。

本项目的链路可以画成：
\begin{center}
\begin{tikzpicture}[node distance=6mm]
\node[flowbox] (a) {AST\\语义结果};
\node[flowbox,right=of a] (b) {结构化 IR\\文本 .ll};
\node[flowbox,right=of b] (c) {Clang 22\\RV32IM 汇编};
\node[flowbox,right=of c] (d) {REIMU\\运行输出};
\draw[flowarrow] (a)--(b);
\draw[flowarrow] (b)--(c);
\draw[flowarrow] (c)--(d);
\end{tikzpicture}
\end{center}

你的生成器负责前两个方框之间的转换。当前语义旁表告诉它“这个调用选中了 \code{PrintlnI32}，这个字面量是 \code{i32} 的 42”；生成器据此创建指令，而不是重新解析 \code{println\_i32} 这个名字。

\subsection{打印 42：完整例子}
配套文件：\code{examples/01\_print42.ll}。其中 \code{println\_i32} 是讲义定义的内部包装函数，输出整数后补一个换行，对应 Rx 的同名内建操作。
\lstinputlisting{examples/01_print42.ll}

\subsection{从文件结构读到执行顺序}
\begin{tabularx}{\linewidth}{@{}p{42mm}Y@{}}
\toprule
代码部分 & 如何理解 \\
\midrule
\code{target datalayout} & 描述目标上的大小和对齐。后面专门解释，初学时先保持项目给出的值。\\
\code{target triple} & 指定 RV32 裸机目标。它是配置，不是运行时指令。\\
\code{@.line} & 一份全局格式字符串。四个字节为 \code{\%}、\code{d}、换行、零字节。\\
\code{declare @printf} & 这里只有接口，具体实现由运行环境提供。\\
\code{define ...} & 定义一个有函数体的函数。\\
\code{entry:} & 基本块的标签；这个名字可以换成 \code{b0}。\\
\code{call ...} & 调用函数。参数写出类型和值。\\
\code{ret ...} & 返回并终结当前基本块。\\
\bottomrule
\end{tabularx}

执行从机器入口 \code{@main} 开始，随后调用 \code{@rx\_source\_main}，再调用本例的 \code{@println\_i32}，最终由 \code{@printf} 输出。源 main 正常结束时返回 \code{void}；最外层机器 main 返回整数 0。

为什么保留两个 main？Rx 的源 \code{main} 是普通的 Unit 函数，源程序也可能递归调用它。机器入口负责把最初那次调用正常完成转换成状态码 0。这个区别与项目的入口规划一致。

\subsection{单独拆开 call 指令}
\begin{lstlisting}
%written = call i32 (ptr, ...) @printf(ptr @.line, i32 %x)
\end{lstlisting}
\begin{itemize}
 \item \code{\%written}：本条指令的结果名，类型为 \code{i32}。
 \item \code{call}：操作码，表示调用。
 \item \code{i32 (ptr, ...)}：被调函数的类型，返回整数，固定参数是一个指针，其余是可变参数。
 \item \code{@printf}：调用目标。
 \item \code{ptr @.line}：格式串地址；\code{i32 \%x}：要打印的整数。
\end{itemize}
printf 的返回值表示打印结果，和源程序要输出的数字不是同一个东西。这里保留了结果名但不再使用它；调用产生的输出副作用仍然发生。

\takeaway{先辨认一段代码是“模块声明”还是“函数内指令”。全局格式串不会按行执行；真正的执行顺序由函数调用和基本块跳转决定。}
\checkq{把源函数中的 42 改成 43，需要改格式字符串吗？答案：不需要，只改变调用实参。}
\srcref{https://releases.llvm.org/22.1.0/docs/LangRef.html\#module-structure}{LLVM 22：模块结构}

\lesson{名字、类型、常量与目标信息}{sec:syntax}
\aim{读懂常见符号和类型写法，并区分“源语言类型”“LLVM 计算类型”“内存存储类型”。}

\subsection{先记住五种写法}
\begin{tabularx}{\linewidth}{@{}p{38mm}Y@{}}
\toprule
写法 & 意义 \\
\midrule
\code{@name} & 模块中的全局对象或函数名字，例如 \code{@main}。\\
\code{\%value} & 函数内 SSA 值、参数；命名结构体类型也使用 \code{\%}，要结合位置判断。\\
\code{entry:} & 定义基本块标签；跳转引用它时写 \code{label \%entry}。\\
\code{; comment} & 分号开始注释，一直到行尾；分号不是语句结束符。\\
\code{\%r = add i32 1, 2} & 定义计算结果 \code{\%r}。这里的等号不会改变已有的 \code{\%r}。\\
\bottomrule
\end{tabularx}

空白主要用来分隔语法元素，缩进用于阅读。建议始终使用有意义的显式名字，例如 \code{\%x.addr}、\code{\%x.value}、\code{loop.test:}，避免初学时同时处理匿名数字编号规则。LLVM 的这些名字和 Rx 局部变量名没有必须一一相同的要求。

\subsection{本项目需要的类型}
\begin{tabularx}{\linewidth}{@{}p{39mm}Y@{}}
\toprule
类型语法 & 读法与用途 \\
\midrule
\code{i1} & 一个比特的整数，通常作为条件值。\\
\code{i8} & 八位整数；项目用于 bool 的内存表示和字符串字节。\\
\code{i32} & 32 位整数。Rx 的四种整数类型都用它承载位模式。\\
\code{ptr} & 指针。当前 LLVM 的指针类型本身不写被指向的元素类型。\\
\code{void} & 函数不返回数据；不是一个可以装进变量的普通值类型。\\
\code{[3 x i32]} & 连续存放的三个 i32。中间的 \code{x} 表示重复数量。\\
\code{\{ i8, i32, i8 \}} & 三个字段组成的结构体类型，通常包含对齐填充。\\
\code{\%Record} & 一个命名类型，如 \code{\%Record = type \{ i8, i32, i8 \}}。\\
\bottomrule
\end{tabularx}

\projectnote{当前规划中，bool 的计算值是 \code{i1}，内存字节是 \code{i8}；数组、结构体和 Vec 的源值由稳定存储承载。这些是生成器的表示选择，不能理解为 LLVM 只允许这一种实现。}

LLVM 的 \code{i32} 没有独立的“有符号类型”和“无符号类型”。同一组 32 位比特如何比较、相除或右移，由所选指令决定。源 TypeId 仍要保存在 lowering 层，供它选择正确的操作。

\subsection{常量和全局数据}
\begin{lstlisting}
@answer = constant i32 42
@counter = global i32 0
@message = private constant [4 x i8] c"%d\0A\00"
\end{lstlisting}
\code{constant} 表示这份全局存储不可修改；\code{global} 表示可修改的全局存储。Rx 当前没有一般全局可变变量需求，第二行只是帮助理解语法。

字符串写法 \code{c"..."} 按字节描述数组。LLVM 的 \code{\textbackslash 0A} 表示十六进制的换行字节，\code{\textbackslash 00} 是 C 字符串终止符；数组长度必须等于实际字节数。这里的 \code{\%d} 位于字符串内部，是 printf 的格式说明，和 SSA 名字无关。

\code{@answer} 作为操作数时表示对象地址。要取得 42，应生成：
\begin{lstlisting}
%value = load i32, ptr @answer
\end{lstlisting}
常用直接常量还包括 \code{i32 42}、\code{i32 -1}、\code{i1 true}、\code{i1 false}、\code{ptr null}。其中 \code{null} 可以表示空指针，但不能被当作有效对象地址解引用。

\subsection{opaque pointer：类型信息转移到了操作上}
\begin{lstlisting}
%p = alloca i32
store i32 42, ptr %p
%v = load i32, ptr %p
\end{lstlisting}
这三行里，\code{\%p} 的 LLVM 类型是 \code{ptr}。load 自己写出“读取 i32”，store 自己写出“写入 i32”。旧资料中可能看到 \code{i32*}；本讲义统一使用 LLVM 22 的 \code{ptr} 形式。

指针不带元素类型，不代表所有读写都合法。如果 \code{\%p} 只指向一个字节，用它读取 i32 依然可能越界或未对齐。项目的 Place 和布局信息负责保留这些约束。

\subsection{目标行如何阅读}
\begin{lstlisting}[numbers=none]
target datalayout = "e-m:e-p:32:32-i64:64-n32-S128"
target triple = "riscv32-unknown-none-elf"
\end{lstlisting}
\begin{tabularx}{\linewidth}{@{}p{28mm}Y@{}}
\toprule
片段 & 本项目中的含义 \\
\midrule
\code{e} & 小端字节序。\\
\code{m:e} & ELF 风格的符号改名规则。\\
\code{p:32:32} & 默认地址空间的指针宽度 32 位、ABI 对齐 32 位。\\
\code{i64:64} & i64 的 ABI 对齐为 64 位，即 8 字节；不是 32 位。\\
\code{n32} & 目标的原生整数宽度集合包含 32 位。\\
\code{S128} & 栈的自然对齐为 128 位，即 16 字节。\\
\bottomrule
\end{tabularx}
\takeaway{datalayout 中这些宽度和对齐数字以比特为单位；指令中的 \code{align 4} 以字节为单位。不要把宿主 Mac 的 64 位指针布局复制到 RV32 模块。}
\srcref{https://releases.llvm.org/22.1.0/docs/OpaquePointers.html}{LLVM 22：Opaque Pointers}

\lesson{函数、参数、调用与返回}{sec:functions}
\aim{写出一个有参数和返回值的函数，理解内部函数、外部声明、Unit 返回和调用签名。}

\subsection{从 Rx 加法函数到 LLVM}
\begin{lstlisting}[language=Rx,numbers=none]
fn add(a: i32, b: i32) -> i32 {
    a + b
}
\end{lstlisting}
对应的 LLVM 函数与调用见 \code{examples/02\_add.ll}：
\lstinputlisting[firstline=13]{examples/02_add.ll}

定义中的 \code{i32 \%a} 声明一个 i32 参数；调用中的 \code{i32 20} 提供一个 i32 实参。\code{ret i32 \%sum} 必须与函数声明的 i32 返回类型相符。函数名与操作码都叫 \code{add} 是允许的：\code{@add} 是符号，指令中的裸 \code{add} 是关键字。

\subsection{declare 和 define}
\begin{lstlisting}
declare i32 @read_number()

define internal void @show(i32 %x) {
entry:
  call void @println_i32(i32 %x)
  ret void
}
\end{lstlisting}
这段是语法片段：若实际调用 \code{@read\_number}，还需要它的实现或外部支持；\code{@show} 所调用的打印函数也需要实现。仅有未使用的外部声明不要求函数立即存在。\code{internal} 表示函数只供当前模块使用；没有它的普通定义可用于导出入口。

前向调用和互递归不要求按调用顺序排列函数体。生成器先登记所有函数签名，再生成各函数体即可；最终对源函数输出定义，对真正外部的函数输出声明。

\subsection{返回 void 的调用}
\begin{lstlisting}
call void @println_i32(i32 42)
ret void
\end{lstlisting}
\code{void} 调用没有可供使用的返回数据，所以不写 \code{\%r = call void ...}。在 Rx 中，\code{println\_i32(...)} 的结果是 \code{()}；在 LLVM 函数签名层面，用 \code{void} 表达无返回数据。

\subsection{可变参数与外部边界}
printf 的第一个固定参数是格式串指针；后面由格式串决定要读取哪些参数。打印 i32 时传 \code{i32}，打印 bool 前应先转换到 i32。单靠 LLVM 类型检查不能证明格式串与可变参数吻合，所以这些调用最好集中在生成器的 I/O 包装中。

实际执行时，LLVM 后端会根据目标调用约定安排参数寄存器、栈和返回值。IR 阶段写出正确的签名、实参顺序和返回类型即可；自行安排 RISC-V 寄存器属于后续 Codegen 工作。

\takeaway{函数调用既产生返回值，也可能产生副作用。不要因为返回值没被使用，就认为调用可以删除；也不要用同一个函数符号承载不一致的签名。}
\checkq{源函数返回 i32，某条路径却写 \code{ret void}，是否只是“忽略结果”？答案：不是，它违反函数返回类型，必须修正这条路径。}
\srcref{https://releases.llvm.org/22.1.0/docs/tutorial/MyFirstLanguageFrontend/LangImpl03.html}{官方教程第 3 章：表达式与函数生成}

\lesson{SSA 值、地址与可变变量}{sec:memory}
\aim{用 alloca、load、store 实现变量更新，并能指出一段代码里的每一个 ptr 指向什么。}

\subsection{三个名字，三种不同含义}
\begin{lstlisting}
%x.addr = alloca i32, align 4
store i32 40, ptr %x.addr, align 4
%old = load i32, ptr %x.addr, align 4
%new = add i32 %old, 2
store i32 %new, ptr %x.addr, align 4
\end{lstlisting}
\begin{center}
\begin{tikzpicture}[node distance=14mm]
\node[flowbox,text width=32mm] (p) {\code{\%x.addr}\\地址值};
\node[flowbox,text width=35mm,right=of p] (mem) {i32 存储槽\\40 改为 42};
\node[flowbox,text width=30mm,right=of mem] (v) {\code{\%old = 40}\\\code{\%new = 42}};
\draw[flowarrow] (p)--node[above,font=\small]{指向}(mem);
\draw[flowarrow] (mem)--node[above,font=\small]{load}(v);
\end{tikzpicture}
\end{center}

\code{\%x.addr} 一直是同一个地址值。\code{\%old} 一直代表那次 load 取得的 40。\code{\%new} 代表加法结果 42。改变的是地址指向的内存内容，而不是这些 SSA 名字已有的定义。

\textbf{SSA（静态单赋值）}要求一个局部结果名在函数中只有一个定义。下面是刻意写错的片段：
\begin{lstlisting}
; INVALID: the same SSA name is defined twice.
%x = add i32 1, 2
%x = add i32 %x, 3
\end{lstlisting}
可改成两个名字，也可通过内存表示源变量：
\begin{lstlisting}
%x1 = add i32 1, 2
%x2 = add i32 %x1, 3
\end{lstlisting}

\subsection{三条内存指令的模板}
\begin{tabularx}{\linewidth}{@{}p{24mm}Y@{}}
\toprule
指令 & 形状与含义 \\
\midrule
alloca & \code{\%p = alloca i32, align 4}：为一个 i32 申请当前调用中的栈存储，结果是 ptr。\\
load & \code{\%v = load i32, ptr \%p, align 4}：从地址读出 i32 值。\\
store & \code{store i32 \%v, ptr \%p, align 4}：把 i32 值写入地址；没有结果值。\\
\bottomrule
\end{tabularx}

\code{alloca} 本身不初始化对象。第一次 load 前必须有适当的写入，且所有可能到达该 load 的执行路径都要满足这一点。LLVM 能解析一个 load，不代表它已经证明该地址保存着正确值。

\code{align 4} 声明此次访问满足至少 4 字节对齐。不要随意把这个数字写大。向一个只保证字节对齐的地址读写 i32 时，必须先明确实际存储和对齐约定。

\projectnote{当前设计把固定 alloca 集中到函数入口的前序区，初始化仍放在真实执行位置。循环回边跳回单独的测试块或循环体，不重复执行入口的 alloca。}

\subsection{可变变量的完整函数}
对应源行为是：令 x 为 40，将 x 加 2，再打印 x。文件 \code{examples/03\_variables.ll} 的主体如下：
\lstinputlisting[firstline=13]{examples/03_variables.ll}
\checkq{最后一次 store 会让 \code{\%old} 变成 42 吗？答案：不会。后续若想读到新内容，需要新的 load。}

\subsection{引用：地址值与保存地址的槽}
令 x 初始为 40，通过 \code{let r = \&mut x;} 取得引用，再执行 \code{*r += 2} 并打印 x。对应完整示例 \code{examples/10\_references.ll}：
\lstinputlisting[firstline=13]{examples/10_references.ll}

逐层看：
\begin{enumerate}
 \item \code{\%x.addr} 指向存放整数的槽。
 \item \code{\%r.addr} 指向存放一个指针值的槽，分配类型为 \code{ptr}。
 \item \code{\%r = load ptr, ptr \%r.addr} 取出指针，结果指向 x。
 \item \code{load i32, ptr \%r} 才读取 x 的整数内容。
\end{enumerate}
虽然两个地址的 LLVM 类型都是 ptr，但它们指向的对象不同。不要因为“都是指针”就少做或多做一次 load。

\takeaway{判断一行代码时，问两件事：操作数现在是地址还是内容？本条指令将产生地址还是内容？这个问题比记忆某个变量名更可靠。}
\srcref{https://releases.llvm.org/22.1.0/docs/tutorial/MyFirstLanguageFrontend/LangImpl07.html}{官方教程第 7 章：可变变量与内存}

\lesson{算术、比较、位运算与转换}{sec:operations}
\aim{根据源整数的有符号性选指令，正确处理 bool、移位和 Rx 的 32 位语义。}

\subsection{二元运算的共同形状}
\begin{lstlisting}
%sum = add i32 %a, %b
%product = mul i32 %sum, 2
\end{lstlisting}
普通二元运算通常写成“结果名 = 操作码 类型 左操作数, 右操作数”。这一行只写一次 \code{i32}，两个操作数都必须匹配它。

\begin{tabularx}{\linewidth}{@{}p{30mm}p{45mm}Y@{}}
\toprule
源操作 & LLVM 指令 & 注意点 \\
\midrule
加、减、乘 & \code{add / sub / mul} & 对 i32 产生 32 位结果。\\
有符号除法、余数 & \code{sdiv / srem} & 除法向零截断。\\
无符号除法、余数 & \code{udiv / urem} & 按无符号位模式解释。\\
按位与、或、异或 & \code{and / or / xor} & bool 的非短路运算也能使用。\\
左移 & \code{shl} & 先按 Rx 规则处理移位计数。\\
有符号右移 & \code{ashr} & 高位补符号位。\\
无符号右移 & \code{lshr} & 高位补零。\\
\bottomrule
\end{tabularx}

没有专门的普通整数“取负”操作码，写 \code{sub i32 0, \%x} 即可。整数按位取反可与全 1 异或；bool 逻辑取反可写：
\begin{lstlisting}
%not.value = xor i32 %x, -1
%not.flag = xor i1 %flag, true
\end{lstlisting}

\subsection{icmp 的结果总是条件值}
\begin{lstlisting}
%same = icmp eq i32 %a, %b
%less.signed = icmp slt i32 %a, %b
%less.unsigned = icmp ult i32 %a, %b
\end{lstlisting}
上面三条结果都是 i1。谓词 \code{eq}、\code{ne} 表示相等、不等；\code{slt}、\code{sle}、\code{sgt}、\code{sge} 表示有符号顺序；\code{ult}、\code{ule}、\code{ugt}、\code{uge} 表示无符号顺序。

例如，位模式全 1 可以作为有符号的 -1，也可以作为无符号的 4294967295。因此 \code{icmp slt i32 -1, 1} 得到 true，\code{icmp ult i32 -1, 1} 得到 false。结果差异来自谓词，而不是 \code{i32} 类型名字。

\subsection{bool 与整数之间的转换}
\begin{lstlisting}
%as.int = zext i1 %flag to i32
%as.byte = zext i1 %flag to i8
store i8 %as.byte, ptr %flag.addr
%stored.byte = load i8, ptr %flag.addr
%loaded.flag = trunc i8 %stored.byte to i1
\end{lstlisting}
\code{zext} 扩展位宽并在高位补零；\code{trunc} 保留低位并缩小位宽。项目约定 bool 的内存字节只为 0 或 1，因此将这个字节截到 i1 可恢复条件值。

\textbf{拓展：}\code{sext} 是符号扩展。把 i1 的 true 符号扩展成 i32 会得到全 1，也就是 -1；所以 Rx 的 \code{true as i32} 应使用 zext，结果为 1。当前项目的核心转换指令先实现 zext/trunc 即可。

Rx 的 i32、u32、isize、usize 都是 32 位。它们之间的 \code{as} 保留位模式，通常不用发出额外的 LLVM 转换指令；生成器仍需更新源类型身份，以影响后续比较、除法和右移。

\subsection{移位计数必须屏蔽}
\begin{lstlisting}
%count = and i32 %raw.count, 31
%left = shl i32 %x, %count
%right.signed = ashr i32 %x, %count
%right.unsigned = lshr i32 %x, %count
\end{lstlisting}
Rx 要求移位计数只使用最低 5 位，因而 33 等价于 1。LLVM 的移位操作不会自动替你履行这个源语言约定；直接将 32 或更大计数用于 i32 移位会产生 poison。普通表达式和复合赋值都要保持同样的处理。

\subsection{属性是承诺，不是装饰}
\begin{lstlisting}
%wrapped = add i32 %a, %b
; Only use this form when signed overflow is proved impossible:
%promised = add nsw i32 %a, %b
\end{lstlisting}
Rx 的普通 32 位加、减、乘和有符号取负按位回绕。基础生成器应使用第一种形式。\code{nsw} 声称没有有符号溢出，\code{nuw} 声称没有无符号溢出；承诺被违反时，LLVM 的含义会改变。不要从其他语言的 Clang 输出中直接复制这些标志。

类似地，第一版无需主动附加 \code{exact}、GEP 的 \code{inbounds} 或引用参数的 \code{noalias}。只有证明满足相应前提，才考虑使用它们。

\takeaway{源语言规则决定选哪条指令、是否需要额外指令。LLVM 能接受 \code{shl} 或 \code{add nsw}，不代表它们自动实现了 Rx 所需的行为。}

\code{examples/12\_scalar\_rules.ll} 把本课的差异串起来，预期输出为：
\begin{lstlisting}[language={},numbers=none]
0
1
0
-4
2147483644
\end{lstlisting}
五行分别来自回绕加法、两种比较、两种右移。你可以先手算，再运行核对。
\srcref{https://github.com/ACMClassCourse-2025/rx-compiler-specification/blob/master/src/expressions/operator-expr.md}{Rx 规范：运算与转换}

\lesson{基本块、分支与短路求值}{sec:branches}
\aim{把 if 和逻辑短路画成控制流图，再逐块写出合法的 LLVM IR。}

\subsection{基本块必须有明确出口}
基本块（basic block）是一段从入口顺序执行、最后由终结指令决定去向的代码。函数体的第一个块是入口块。块的文本排列顺序不决定执行顺序：LLVM 没有“这一块写完，就自动落到下一块”的规则。

\begin{lstlisting}
br label %next
br i1 %condition, label %yes, label %no
ret i32 %value
ret void
\end{lstlisting}
以上是四种不同的出口模板，不应连续放进同一个块。无条件 br 直接跳转；条件 br 需要 i1，true 选择第一个目标，false 选择第二个目标；ret 结束当前函数调用。

\takeaway{每个基本块必须以一条终结指令结束。终结指令后不能继续往该块追加指令。要执行下一段代码，就创建或选择相应的基本块。}

\subsection{if 表达式：用槽保存合流结果}
下面的源函数根据条件选择一个值：
\begin{lstlisting}[language=Rx,numbers=none]
fn max(a: i32, b: i32) -> i32 {
    if a > b { a } else { b }
}
\end{lstlisting}
控制流图中的箭头表示可能的跳转：
\begin{center}
\begin{tikzpicture}[node distance=8mm and 22mm]
\node[flowbox] (entry) {entry\\比较 a 与 b};
\node[flowbox,below left=of entry] (yes) {then\\保存 a};
\node[flowbox,below right=of entry] (no) {else\\保存 b};
\node[flowbox,below right=of yes] (merge) {merge\\读出结果并返回};
\draw[flowarrow] (entry)--node[above left,font=\small]{true}(yes);
\draw[flowarrow] (entry)--node[above right,font=\small]{false}(no);
\draw[flowarrow] (yes)--(merge);
\draw[flowarrow] (no)--(merge);
\end{tikzpicture}
\end{center}
配套文件 \code{examples/04\_if\_slots.ll} 中的函数如下。外层 main 调用 \code{max\_slot(17, 42)}，输出 42。
\lstinputlisting[firstline=13,lastline=28]{examples/04_if_slots.ll}

result.addr 在入口创建，因此两个分支都能访问。同一次调用只会走进其中一个分支，而该分支会先写结果槽再跳到 merge。merge 中读取的是本次实际执行路径留下的值。

如果某个分支已经 ret，它不会跳到 merge，也不必给 merge 的结果槽写值。生成器需要检查的是：\textbf{所有正常抵达合流块的路径都提供了结果}。若两个分支都不会正常结束，整个 if 没有正常结果，不应再读取结果槽。

\subsection{短路意味着右边可能完全不执行}
Rx 的 \code{\&\&} 和 \code{||} 会根据左操作数决定是否求右操作数。不能先把左右两边全部生成并执行，再用 \code{and} 或 \code{or} 合并结果。

配套文件 \code{examples/05\_short\_circuit.ll} 中，\code{@rhs} 会打印 999，然后返回 true。下面的函数实现 \code{left \&\& rhs()}：
\lstinputlisting[firstline=19,lastline=35]{examples/05_short_circuit.ll}

\begin{itemize}
 \item left 为 false：直接进入 skip.rhs，保存 false，\code{@rhs} 不会被调用。
 \item left 为 true：进入 eval.rhs，调用右侧并保存其结果。
 \item 两条路径在 merge 汇合；bool 仍遵守“计算 i1、存储 i8”的项目约定。
\end{itemize}
main 先调用 false，再调用 true，并把两次逻辑结果转为整数打印。完整程序输出为：
\begin{lstlisting}[language={},numbers=none]
0
999
1
\end{lstlisting}
第一行前没有 999，正是短路生效的证据。实现 \code{left || rhs()} 时，left 为 true 的路径直接保存 true；只有 false 路径才调用 rhs。

\textbf{拓展：}\code{select} 可以根据条件选择两个已有 SSA 值，但不会撤销为计算它们已经发生的调用、读写或除法。因此它不能直接代替含副作用表达式的短路控制流；当前项目也不要求先实现它。

\subsection{unreachable 不是“什么也不做”}
\begin{lstlisting}
dead:
  unreachable
\end{lstlisting}
这条终结指令声明执行不应抵达此处。真的执行到它属于未定义行为。它不是普通 return，也不保证触发一个友好的错误提示，更不能拿来替代本应正常运行的分支。

\projectnote{当前生成器可以保留一个没有正常前驱的预建 merge/exit 块，并用 unreachable 结束它。这依赖于确实没有有效执行路径抵达该块。Unit 表示正常结束但没有数据；NoReturn 表示没有正常后继，不能混成一种情况。}

\checkq{如果将例子中的 \code{call @rhs} 移到 entry、条件跳转之前，false 路径还会短路吗？答案：不会，输出会多出 999。}
\srcref{https://releases.llvm.org/22.1.0/docs/LangRef.html\#terminator-instructions}{LLVM 22：终结指令}

\lesson{循环、支配关系与可选的 phi}{sec:loops}
\aim{写出能更新局部变量的循环，理解跨基本块使用 SSA 值的限制。}

\subsection{while 是一条带回边的路径}
考虑从 1 累加到 n 的函数。这里只讨论不会使循环计数溢出的输入，例如 0 和 10：
\begin{lstlisting}[language=Rx,numbers=none]
fn sum_to(n: i32) -> i32 {
    let mut i = 1;
    let mut sum = 0;
    while i <= n {
        sum += i;
        i += 1;
    }
    sum
}
\end{lstlisting}
可以把它拆成 entry、test、body、exit 四块：
\begin{center}
\begin{tikzpicture}[node distance=8mm and 22mm]
\node[flowbox] (entry) {entry\\创建槽并初始化};
\node[flowbox,below=of entry] (test) {test\\检查 i $\leq$ n};
\node[flowbox,below=of test] (body) {body\\累加并更新 i};
\node[flowbox,right=of test] (exit) {exit\\读取 sum 并返回};
\draw[flowarrow] (entry)--(test);
\draw[flowarrow] (test)--node[right,font=\small]{true}(body);
\draw[flowarrow] (test)--node[above,font=\small]{false}(exit);
\draw[flowarrow] (body.west)--++(-12mm,0)|-(test.west);
\end{tikzpicture}
\end{center}
\lstinputlisting[firstline=13,lastline=34]{examples/06_loop.ll}

每次回到 test 都会重新 load i，所以能读到上一次循环体写入的新值。文本中 \code{\%i} 只有一处定义；这条 load 在运行时执行多次并不违反“静态单赋值”。SSA 的“静态”指的是程序结构里的定义位置，而不是指令一生只能执行一次。

\begin{itemize}
 \item while 的 continue 跳到 test，使条件重新求值。
 \item while 的 break 跳到 exit。
 \item 嵌套循环应使用当前内层循环的目标；外层目标不能被覆盖后丢失。
 \item 入口 alloca 不在回边上，避免每次迭代重新申请局部存储。
\end{itemize}
\code{examples/06\_loop.ll} 依次测试 n=10 和 n=0，输出 55 和 0。第二个输入能检查“第一次条件就为 false”的路径。

\subsection{定义必须支配普通使用}
先用自然语言理解支配（dominance）：如果从函数入口走到使用位置的每一条路径，都必须经过某个定义，那么这个定义支配该使用。普通 SSA 操作数的定义必须满足这一要求；同块内还要先定义再使用。

下面是错误片段：
\begin{lstlisting}
; INVALID: %only.then is absent on the else path.
then:
  %only.then = add i32 %a, 1
  br label %merge
else:
  br label %merge
merge:
  ret i32 %only.then
\end{lstlisting}
如果从 else 到达 merge，就没有经过 \code{\%only.then} 的定义。名字在文本中出现得更早，不能替代控制流上的保证。

本项目第一版可以用结果槽解决这类问题：在入口分配槽，每条正常前驱路径分别保存值，在 merge 读取。槽方案仍需保证“每条抵达路径都初始化”，LLVM 的一般类型检查不会自动证明你的源变量一定赋过值。

\clearpage
\subsection{可选阅读：phi 在合流处选择来源}
LLVM 还提供 phi，直接把不同前驱边带来的值合并成一个 SSA 结果。下面与第 6 课的 max 函数行为相同：
\lstinputlisting[firstline=13,lastline=24]{examples/09_phi.ll}

核心一行按“值、前驱块”成对阅读：
\begin{lstlisting}
%result = phi i32 [ %a, %then ], [ %b, %else ]
\end{lstlisting}
刚才从 then 进入 merge，结果就取 a；从 else 进入，结果就取 b。phi 不重新比较条件，也不会重新执行两个分支。

写 phi 时记住：
\begin{enumerate}
 \item phi 必须位于所在基本块的开头，排在普通指令之前。
 \item incoming 项要匹配实际前驱；不能遗漏某条可能进入的前驱边。
 \item 每个 incoming 值必须在对应前驱边上可用，而不要求它支配合流块内的所有路径。
 \item 所有 incoming 值都要匹配 phi 声明的结果类型。
\end{enumerate}
循环也可以在头部用 phi 选择“第一次进入时的初值”和“回边带来的下一次值”。这留待理解 SSA 构造或优化时再做。

\projectnote{当前核心 Opcode 没有 Phi，分支结果采用槽表示。学习 phi 能帮助阅读 LLVM 教程和优化结果，但实现第一个可用 IR 生成器不必先做 SSA 重命名、支配树或 mem2reg。}
\takeaway{栈槽保存可变状态，phi 合并沿控制流传来的 SSA 值。两种机制都能表达分支和循环，基础实现应先把自己的表示约定贯彻一致。}
\srcref{https://releases.llvm.org/22.1.0/docs/tutorial/MyFirstLanguageFrontend/LangImpl05.html}{官方教程第 5 章：控制流与 phi}

\lesson{数组与 getelementptr}{sec:gep}
\aim{读懂 GEP 的每一个索引，能区分“算地址”和“读取元素”。}

\subsection{数组是连续的元素存储}
\begin{lstlisting}
%items = alloca [3 x i32], align 4
%item.1 = getelementptr [3 x i32], ptr %items, i32 0, i32 1
store i32 42, ptr %item.1, align 4
%answer = load i32, ptr %item.1, align 4
\end{lstlisting}
第一行分配整个数组，结果 \code{\%items} 是它的地址。第二行只计算索引为 1 的元素地址，第三行才写入 42，第四行才读回内容。

\begin{center}
\begin{tikzpicture}[font=\small]
\node[anchor=east] at (-0.15,0.35) {索引};
\node[anchor=east] at (-0.15,-0.4) {字节偏移};
\foreach \x/\i/\off in {0/0/0,3/1/4,6/2/8} {
  \draw[fill=Pale,draw=Rule] (\x,0) rectangle ++(3,0.7);
  \node at (\x+1.5,0.35) {\code{i32} / \i};
  \node[text=Teal] at (\x+1.5,-0.4) {\off};
}
\node[anchor=west,text=Muted] at (9.15,0.35) {总大小 12 字节};
\end{tikzpicture}
\end{center}

\subsection{两个索引为什么是 0 和 1}
把这一行拆开看：
\begin{lstlisting}
%item.1 = getelementptr [3 x i32], ptr %items, i32 0, i32 1
\end{lstlisting}
\begin{tabularx}{\linewidth}{@{}p{40mm}Y@{}}
\toprule
部分 & 意义 \\
\midrule
\code{[3 x i32]} & 基础元素的布局类型：从一个“完整数组”出发计算。\\
\code{ptr \%items} & 基地址。\\
第一个 \code{i32 0} & 在完整数组之间移动 0 份，仍然是当前这个数组。\\
第二个 \code{i32 1} & 进入数组内部，选择索引为 1 的 i32 元素。\\
结果 \code{\%item.1} & ptr 类型的元素地址，不是元素的 i32 值。\\
\bottomrule
\end{tabularx}
在这个例子里，两个索引对应的地址偏移可理解为
\[
0 \times \operatorname{sizeof}([3 \times i32])
+ 1 \times \operatorname{sizeof}(i32)=4\text{ 字节}.
\]
第一个索引是指针层面的步进，不是一次 load，也不是 LLVM 自动补上的“解引用运算”。

如果手里已经是连续 i32 数据的起始地址，写法更短：
\begin{lstlisting}
%element = getelementptr i32, ptr %data, i32 %index
\end{lstlisting}
这里基础类型是 i32，因此一个索引就按 4 字节步进。不要机械地给每一个 GEP 都添加 \code{0, index}；索引形状要由基础对象类型决定。

\subsection{完整数组示例}
\code{examples/07\_arrays.ll} 初始化三个元素，再用运行时计算的索引读取第二个：
\lstinputlisting[firstline=13]{examples/07_arrays.ll}
输出 42。把 \code{\%index = add i32 0, 1} 改为加 2，应该输出 30；这是检查索引选择是否正确的简单实验。

\subsection{嵌套数组与运行时索引}
对 \code{[2 x [3 x i32]]}，一条 GEP 可以逐层进入：
\begin{lstlisting}
%cell = getelementptr [2 x [3 x i32]], ptr %matrix,
                     i32 0, i32 %row, i32 %column
\end{lstlisting}
这是一条分行展示的指令：先停留在当前矩阵，再选择行数组，最后选择行中的元素。数组索引允许是 SSA 值；结构体字段索引则有更严格的限制，见下一课。

\textbf{GEP 不做边界检查。}合法指针运算和合法内存访问是两件事。算出一个地址不表示可以读取它。若源语言要求边界行为，应另外实现对应逻辑；当前项目按规范中有效输入的约束处理，不能把 GEP 当成自带检查的数组 API。

\code{getelementptr inbounds} 还附带额外的指针运算约束，违反时可能产生 poison。初学及基础生成器可先不添加这个承诺。即使省略 inbounds，也不能因此合法地访问不存在的对象。

\takeaway{GEP 回答“目标在哪里”，load 回答“目标存着什么”。索引的每一步由显式的布局类型解释，opaque ptr 不会替你保存这些信息。}
\checkq{读取 \code{a[1]} 为什么常有 GEP 和 load 两步？答案：先从数组地址计算元素地址，再从该地址读值；只取 \code{\&a[1]} 时可以停在地址这一步。}
\srcref{https://releases.llvm.org/22.1.0/docs/GetElementPtr.html}{LLVM 22：The Often Misunderstood GEP Instruction}

\lesson{结构体、字段布局与聚合返回}{sec:structs}
\aim{用字段序号寻址，理解填充字节，以及项目通过地址传递聚合值的方式。}

\subsection{命名结构体与字段 GEP}
\begin{lstlisting}
%Record = type { i8, i32, i8 }

; Inside a function:
%record = alloca %Record, align 4
%value.addr = getelementptr %Record, ptr %record, i32 0, i32 1
store i32 42, ptr %value.addr, align 4
\end{lstlisting}
\code{\%Record} 出现在 type 定义处，是类型名字；\code{\%record} 是函数里的值名字。字段索引从 0 开始，因此 1 选中第二个字段 i32。

LLVM 不认识源字段名字，例如 \code{record.value} 中的 value。semantic 已经确认了字段归属，lowering 应使用声明顺序中的字段号，而不是按字段名字排序。

结构体的字段索引必须是常量 \code{i32}，不能用一个运行时 SSA 整数任意选择不同类型的字段。数组里的元素类型相同，所以可以使用运行时下标；结构体各字段的类型和偏移可能不同。

\subsection{三个字段为什么占 12 字节}
按本项目的 RV32 布局，\code{\{ i8, i32, i8 \}} 的布局为：

\begin{tabularx}{\linewidth}{@{}p{24mm}p{30mm}p{33mm}Y@{}}
\toprule
区域 & 字节偏移 & 大小 & 原因 \\
\midrule
字段 0 & 0 & 1 字节 & i8 按 1 字节对齐。\\
填充 & 1--3 & 3 字节 & 下一个 i32 需要按 4 字节对齐。\\
字段 1 & 4 & 4 字节 & i32 占 4 字节。\\
字段 2 & 8 & 1 字节 & 紧接前一个字段。\\
尾部填充 & 9--11 & 3 字节 & 整体大小向最大字段对齐 4 取整。\\
\bottomrule
\end{tabularx}

总大小是 12 字节，整体对齐是 4 字节。如果创建一个该结构体的数组，相邻元素的起始地址相差 12 字节，而不是字段大小之和 6。

GEP 根据 datalayout 和结构体类型完成偏移计算，不要求把 4 或 8 手工写成字节偏移。项目自己的布局计算仍需与之保持一致，因为分配、表示复制和后续自写后端也会依赖这些数字。

LLVM 还有写成 \code{<\{ i8, i32, i8 \}>} 的 packed 结构体，它改变布局规则；本项目的普通结构体不应为了“省空间”随意改用这种形式。

\subsection{bool 字段的完整例子}
\code{examples/08\_structs.ll} 的主体如下，输出 42 和 1：
\lstinputlisting[firstline=13]{examples/08_structs.ll}
读取 bool 字段时先 load i8，再 trunc 到 i1；打印前再 zext 到 i32。三种位宽在不同边界上各有用途，不能因为最终打印的是整数就跳过表示约定。

\subsection{聚合地址不自动等于独立的值}
LLVM 本身允许数组或结构体作为某些指令、参数和返回的值。本项目选择另一套便于实现的内部约定：数组、结构体和 Vec 的值由稳定存储承载，ABI 中传递其地址。

这里要区分两个指针：
\begin{itemize}
 \item 指向源可变对象的地址：对象后续可能被修改。
 \item 指向本次求值快照的地址：内容在这次值使用期间必须保持稳定。
\end{itemize}
两者的 LLVM 类型都是 ptr，但不能因为类型相同就互换。源函数按值接收数组时，不能意外变成对调用者数组的可变别名。

\subsection{通过隐藏输出地址返回数组}
配套文件 \code{examples/11\_aggregate\_return.ll} 用下面的形状返回两个整数：
\lstinputlisting[firstline=13]{examples/11_aggregate_return.ll}

可以按三步读这个例子：
\begin{enumerate}
 \item 调用者分配 \code{[2 x i32]} 的结果存储。
 \item 把地址作为额外参数传给 \code{@make\_pair}。
 \item 被调方往该地址写入 20 和 22，以 \code{ret void} 结束；调用者读取并相加，得到 42。
\end{enumerate}
这个输出参数是双方约定的普通 ptr。仅仅把它放在参数列表第一位，就能实现本例，不必先学习或随意添加 LLVM 的 \code{sret} 属性。

\projectnote{当前项目约定聚合返回使用隐藏的首个输出地址；聚合参数传入值快照地址，被调方复制到自己的参数槽。复制对象表示不等于显式 Clone：Box、Vec 的深复制有不同的源语义，不能混用。}
\takeaway{LLVM 类型负责说明机器对象的形状；lowering 还必须保证源程序要求的独立值、求值顺序和存储生命周期。只检查 ptr 类型一致远远不够。}

\lesson{把语法接到当前 Rx 编译器}{sec:project}
\aim{从现有 semantic 结果选择指令，建立正确的值、位置和函数边界，而不是重新做语义分析。}

\subsection{先把源概念对应到 IR}
\begin{tabularx}{\linewidth}{@{}p{31mm}p{45mm}Y@{}}
\toprule
源概念 & 常见 IR 表示 & lowering 仍需知道什么 \\
\midrule
四种整数 & i32 & 有符号性、运算规则。\\
bool & 计算 i1，存储 i8 & 在 load/store 边界转换。\\
引用、Box & 标量 ptr & 指向的源类型及具体操作。\\
数组、结构体 & 指向稳定存储的 ptr & 对象布局、复制和字段顺序。\\
Vec & 指向三字段对象的 ptr & 元素类型、长度、容量与数据区。\\
Unit & 正常结果，无数据 & 能继续生成后续指令。\\
never / NoReturn & 无正常结果 & 停止沿当前路径生成。\\
\bottomrule
\end{tabularx}

TypeId 的身份不能混用。semantic 的 TypeId 区分源类型；核心 IR 的 TypeId 描述 LLVM 层的类型。源 i32 与 u32 可以对应同一个 IR i32，但 sdiv/udiv 的选择仍须来自源类型。

\subsection{Place、读取与取地址}
考虑三个很相似的源表达式：
\begin{lstlisting}[language=Rx,numbers=none]
x
&x
*r
\end{lstlisting}
当 x 处在需要值的位置，生成器从 x 的槽读取；\code{\&x} 取得槽地址；\code{*r} 先取得引用的 ptr，再根据本次使用场景决定是保留目标位置，还是读取目标内容。

semantic 的 Place 是已经检查过的求址计划；IR lowering 的 Place 则含运行时地址及存储信息。将字段投影翻译成 GEP，将读取翻译成 load，将写入翻译成 store，就把前端计划连接到了具体指令。

\subsection{普通指令也可能组合出错误的源行为}
赋值不仅是“最后生成一条 store”。当前 Rx 项目要求普通赋值和复合赋值先求右侧，再求左侧目标地址；复合赋值随后读取目标旧值、计算并写回。地址表达式只能按要求求值一次。

同样，调用实参按源顺序求值，并立即保存需要的快照。如果第一个实参是数组值，后一个实参会修改这个数组，不能拖到所有实参求完后才复制第一个参数。

二元运算的引用操作数还涉及“何时读取所指目标”。这类规则应以 \code{docs/ir.md} 的求值顺序约定为准，不能只根据 LLVM 的操作数列表猜测。先写出预期副作用顺序，再决定每次 load 放在哪个位置。

\subsection{Vec：先取字段，再访问数据}
当前项目把 Vec 的存储布局表示成：
\begin{lstlisting}
%VecI32 = type { ptr, i32, i32 }
\end{lstlisting}
字段依次为 data、len、cap。在 RV32 中三个字段均占 4 字节，因此整个对象为 12 字节、按 4 字节对齐。本机 64 位平台不能据此照抄这个数字。

在已确认下标有效的前提下，元素寻址片段为：
\begin{lstlisting}
%data.addr = getelementptr %VecI32, ptr %vec, i32 0, i32 0
%data = load ptr, ptr %data.addr
%element.addr = getelementptr i32, ptr %data, i32 %index
%element = load i32, ptr %element.addr
\end{lstlisting}
第一个 GEP 只找到“保存 data 指针的字段”。必须 load 出真正的数据区地址后，才能对元素做下一次 GEP。GEP 不会自动穿过一个存储在字段里的指针。

Box 的分配、Vec 扩容和 Clone 最终也会由调用、寻址、读写与循环组合而成。学习语法时先掌握这些基础，运行时 helper 的完整实现仍按项目规划推进。

\projectnote{不能用 \code{icmp eq ptr} 自动实现所有源 equality。LLVM 的这种比较检查指针值；当前项目中引用、Box 的相等操作要比较所指内容，Vec 要比较长度及有效元素，结构体不能把 padding 也拿来整体比较。操作选择必须来自源类型语义。}

\subsection{从例子映射到 Builder}
目前 \code{include/ir/ir.hpp} 中的 Builder 接口，已经能与本讲义的基础语法对应：

\begin{tabularx}{\linewidth}{@{}p{43mm}Y@{}}
\toprule
接口概念 & 最终文本 \\
\midrule
\code{slot / load / store} & alloca、load、store。\\
\code{gep} & 带基础存储类型和索引列表的 getelementptr。\\
\code{binary / compare} & 算术或位运算、带谓词的 icmp。\\
\code{convert} & zext 或 trunc 等转换。\\
\code{call} & 与函数签名匹配的 call。\\
\code{jump / branch / ret} & 无条件 br、条件 br、ret。\\
\bottomrule
\end{tabularx}

先创建结构化 Module、Function、Block 和指令，再验证并打印 \code{.ll}。printer 只负责把已有信息写成正确语法，不应重新查 AST 决定类型，也不应通过补默认返回掩盖缺失的控制流。

\textbf{第一个接入目标。}对源 \code{fn main() \{ println\_i32(42); \}}，从已解析的调用目标识别内建打印，生成第 1 课的内部源函数与机器入口。等这个链路能输出 42，再逐步增加本讲义后面的指令组合。这里给出接入思路；配套示例由手写 IR 构成，并不代表项目 CLI 已完成这些功能。

\lesson{验证、运行与排错}{sec:workflow}
\aim{建立“手写小例子、检查 IR、编译、运行、核对输出”的学习闭环。}

\subsection{三种检查回答不同问题}
\begin{tabularx}{\linewidth}{@{}p{35mm}Y@{}}
\toprule
检查层次 & 能回答的问题 \\
\midrule
IR 结构验证 & 类型、SSA 使用、基本块终结和指令约束是否成立。\\
本机执行 & 在宿主平台上，例子的输出是否符合预期。\\
RV32 目标执行 & 目标布局、调用约定、RISC-V 后端与 REIMU 是否一起工作。\\
\bottomrule
\end{tabularx}
第三层是项目验收必需的。前两层有助于学习和定位问题，但不能替代目标链路。即使三层都通过一组例子，也只验证这组例子覆盖到的行为。

\subsection{先在本机跑完配套例子}
以下命令从仓库根目录开始，只把产物写进 \code{docs/LLVM}：
\begin{lstlisting}[language={},numbers=none]
cd docs/LLVM
python3 verify_examples.py --mode native --clang /usr/bin/clang
\end{lstlisting}
脚本读取 \code{examples/manifest.json}，在 \code{build/native/} 生成去掉两行目标信息的副本，分别编译、运行，并逐字节比对 \code{.out} 文件。原始 RV32 示例保持原样。

交付时的实际检查结果是：\textbf{12 个示例全部通过本机编译和输出核对}。使用 Apple Clang 21，在 ARM64 macOS 执行。本次环境未提供 \code{clang-22}，因此没有宣称 RV32 / REIMU 目标运行已通过。编译器版本、实际命令与逐例输出保存在 \code{build/validation-report.native.json}。

\subsection{使用 LLVM 22 做结构检查}
准备好带版本号的 LLVM 22 工具后，可在 \code{docs/LLVM} 中执行：
\begin{lstlisting}[language={},numbers=none]
mkdir -p build/check
llvm-as-22 examples/01_print42.ll -o build/check/01_print42.bc
opt-22 -passes=verify -disable-output build/check/01_print42.bc
\end{lstlisting}
\code{llvm-as} 把可读文本转换成 bitcode；\code{opt} 的 verify pass 检查 IR 的有效性。不同安装方式可能把工具命名为 \code{llvm-as}、\code{opt}，应先核对版本再替换命令名。

这不会执行程序，也不验证 printf 的格式串含义、源语言求值顺序或数组快照是否正确。

\subsection{项目目标链路：LLVM 22 到 RV32IM}
在 \code{docs/LLVM} 中，单个示例可按项目同样的目标参数编译：
\begin{lstlisting}[language={},numbers=none]
mkdir -p build/target
clang-22 -S -x ir examples/01_print42.ll \
  -o build/target/01_print42.s \
  --target=riscv32-unknown-none-elf \
  -march=rv32im -mabi=ilp32 -O0 \
  -mllvm -riscv-no-aliases
python3 ../../scripts/strip_asm_debug.py build/target/01_print42.s
\end{lstlisting}
\code{-S} 输出汇编，\code{-x ir} 明确输入语言，target/march/mabi 指定 RV32IM / ILP32；其余参数与项目运行链路一致。注意：不能用本机 ARM64 汇编代替这里的 RISC-V 汇编。

已有 REIMU 可执行文件时，可让配套脚本批量执行完整目标链路：
\begin{lstlisting}[language={},numbers=none]
python3 verify_examples.py --mode target \
  --clang /absolute/path/to/clang-22 \
  --reimu /absolute/path/to/reimu
\end{lstlisting}
两处绝对路径需要替换为自己的工具位置。脚本要求 Clang 22，输出汇编、日志和模拟器结果到 \code{build/target/}；缺少工具时明确失败。报告中的
\code{target\_\allowbreak execution\_\allowbreak verified}
只有在目标模式全部通过时才会为 true。

\subsection{自己改例子时的顺序}
\begin{enumerate}
 \item 复制一个已通过的例子到本目录的 \code{build/experiments/}。
 \item 一次只改变一个因素：常量、谓词、字段号或一条控制流边。
 \item 先手算输出，再编译运行。
 \item 观察错误属于解析、验证、编译还是运行输出不符。
 \item 最后才把相同结构搬进自己的 IR Builder。
\end{enumerate}
大多数例子直接用返回 i32 的机器 main 来聚焦语法；第 1 个例子额外展示内部源 main 包装。将它们接到 Rx 编译器时，应恢复项目统一的源函数与机器入口约定。

\subsection{常见错误的定位表}
\begin{tabularx}{\linewidth}{@{}p{47mm}Y@{}}
\toprule
现象 & 优先检查 \\
\midrule
结果名重复、解析失败 & 是否对同一 \code{\%name} 定义两次；每个新计算是否获得新名字。\\
operand 类型不匹配 & i1/i8/i32 是否在边界转换；load 写的是内容类型还是指针类型。\\
基本块缺少终结指令 & 是否漏掉 ret/br；是否以为标签之间会自动顺序执行。\\
定义不支配使用 & 值是否只在一个分支中定义，却被所有合流路径使用。\\
字符串数组类型不匹配 & 数组长度是否包含换行和末尾零字节。\\
输出了额外内容 & 短路右侧是否提前调用；循环退出后是否仍生成了副作用。\\
结构体字段读错 & 字段号、声明顺序、存储类型和目标布局是否一致。\\
结果偶尔不稳定 & 是否读了未初始化槽、过期地址，或写入了错误对象。\\
本机成功、目标失败 & 编译器版本、target/march/mabi、指针宽度和运行时接口。\\
\bottomrule
\end{tabularx}

\lesson{动手练习与参考答案}{sec:practice}
\aim{不用背诵大段语法，能够预测结果、修改示例，并解释修改为什么成立。}

\subsection{练习：先独立回答}
\begin{enumerate}[label=\textbf{\arabic*.}]
 \item \textbf{从 42 到 43。}修改第 1 个例子，让它打印 43。需要改变格式串、函数签名或机器 main 的返回值吗？
 \item \textbf{旧值与新值。}在第 3 个例子最后一次 store 后，依次打印 \code{\%old} 和重新 load 得到的值。两行分别是什么？
 \item \textbf{有符号比较。}让 max 的实参变成 -1 和 1。分别使用 sgt 和 ugt 比较，会选出哪个值？若都以 \code{\%d} 打印，屏幕上是什么？
 \item \textbf{合流槽。}删除第 4 个例子 else 块中的 store，但保留 br。语法上可写出来，为什么仍然不是正确的 max？
 \item \textbf{短路副作用。}第 5 个例子调用 false、true 时输出什么？若把 rhs 调用提前到条件跳转前，输出会如何变化？
 \item \textbf{循环目标。}第 6 个例子中 continue 应跳向哪个块？为什么不能跳回 entry？n 为 0 时是否进入 body？
 \item \textbf{GEP 的步长。}对 \code{[3 x i32]} 的地址，\code{0, 2} 选择哪里？\code{1, 0} 和它相同吗？
 \item \textbf{字段布局。}RV32 中 \code{\{ i8, i32, i8 \}} 的第二个字段偏移和整体大小分别是多少？如果仅把三个字段大小相加，会漏掉什么？
 \item \textbf{bool 转换。}将 i8 的数值 2 trunc 到 i1 得到什么？若要测试任意 i32 是否非零，应使用什么指令？
 \item \textbf{返回聚合。}为何第 11 个例子让调用者分配结果槽，再把地址传给被调方？能否改成返回被调方局部 alloca 的地址，并在返回后读取？
\end{enumerate}

\subsection{参考答案}
\begin{enumerate}[label=\textbf{\arabic*.}]
 \item 将源函数中传给打印包装的实参改成 \code{i32 43}。格式串仍打印十进制整数，签名不变，机器 main 仍返回状态码 0。
 \item 依次为 40、42。\code{\%old} 保存先前那一次 load 的结果；store 改的是内存，不会回头修改已有 SSA 值。
 \item sgt 按有符号解释，选择 1；ugt 把 -1 的位模式看作最大的 32 位无符号整数，选择全 1 的位模式。用 \code{\%d} 打印时分别看到 1 和 -1。比较方式与打印解释是两个独立边界。
 \item 当 a 不大于 b，路径会从 else 到 merge，却没有给结果槽写入 b。merge 的 load 不能得到定义良好的 max 结果。补回 store，或采用正确的 phi 合流。
 \item 正确顺序是 0、999、1。提前调用后，false 那次也打印 999，于是变成 999、0、999、1。结果布尔值可能看似相同，副作用已违反短路语义。
 \item continue 跳到 test。entry 包含分配与初始赋值，跳回去会重置状态，而且函数入口块不能有前驱。n=0 时第一次检查失败，直接到 exit，结果为 0。
 \item \code{0, 2} 从当前数组内部取第 3 个元素，偏移 8 字节。\code{1, 0} 先跨过整个 12 字节数组，再取下一份数组的第一个元素；若只分配了这一份数组，不能通过该地址读取另一个数组。
 \item 偏移 4，总大小 12。字段之间与结构体尾部都可能有 padding，不能只算 1+4+1。
 \item 2 的最低位为 0，所以 trunc 得到 false。测试一般整数是否非零应写 \code{icmp ne i32 \%x, 0}；trunc 不是一般意义的“非零转 true”。
 \item 调用者的槽在被调函数返回后仍有效。被调方局部 alloca 的存储在该调用结束后不再有效，返回其地址再读取不成立。隐藏输出地址使结果从一开始就写到调用者提供的存储中。
\end{enumerate}

\subsection{完成入门的自检标准}
能够在不看答案的情况下完成以下四件事，就可以把主要精力转回项目实现：
\begin{itemize}
 \item 读一段 IR，给每个 SSA 值标明类型，并说明每个 ptr 指向的对象。
 \item 画出分支和循环的基本块，指出每个块的终结指令与可能前驱。
 \item 手写变量更新、短路、数组索引与字段访问，并预测输出。
 \item 对照 semantic 结果解释“为什么选这条指令”，尤其是有符号运算、bool 存储和聚合快照。
\end{itemize}
不必等到掌握全部 LLVM 指令才开始编写 IR。以第 1 个例子为最小闭环，每实现一类源表达式，就补一个能区分正确与错误行为的小程序。

\appendix
\lesson{常用语法速查}{app:syntax}
本表是片段模板，不是完整模块。假定使用前相应的参数、地址、标签和外部函数已经定义，且类型与存储有效。

{\small
\begin{longtable}{@{}p{28mm}p{105mm}p{25mm}@{}}
\toprule
用途 & 语法形状 & 结果类型 \\
\midrule
\endfirsthead
\toprule
用途 & 语法形状 & 结果类型 \\
\midrule
\endhead
\bottomrule
\endfoot
整数加法 & \code{\%r = add i32 \%a, \%b} & i32 \\
整数取负 & \code{\%r = sub i32 0, \%x} & i32 \\
有符号除法 & \code{\%r = sdiv i32 \%a, \%b} & i32 \\
无符号余数 & \code{\%r = urem i32 \%a, \%b} & i32 \\
按位与 & \code{\%r = and i32 \%a, \%b} & i32 \\
bool 取反 & \code{\%r = xor i1 \%flag, true} & i1 \\
算术右移 & \code{\%r = ashr i32 \%x, \%count} & i32 \\
逻辑右移 & \code{\%r = lshr i32 \%x, \%count} & i32 \\
相等比较 & \code{\%r = icmp eq i32 \%a, \%b} & i1 \\
有符号小于 & \code{\%r = icmp slt i32 \%a, \%b} & i1 \\
无符号小于 & \code{\%r = icmp ult i32 \%a, \%b} & i1 \\
bool 扩展 & \code{\%r = zext i1 \%flag to i32} & i32 \\
字节转条件 & \code{\%r = trunc i8 \%byte to i1} & i1 \\
栈槽 & \code{\%p = alloca i32, align 4} & ptr \\
读整数 & \code{\%v = load i32, ptr \%p, align 4} & i32 \\
写整数 & \code{store i32 \%v, ptr \%p, align 4} & 无 \\
读指针字段 & \code{\%p = load ptr, ptr \%slot} & ptr \\
数组元素地址 & \code{\%p = getelementptr [3 x i32],}\newline
  \code{ptr \%a, i32 0, i32 \%i} & ptr \\
结构体字段 & \code{\%p = getelementptr \%Record,}\newline
  \code{ptr \%r, i32 0, i32 1} & ptr \\
函数调用 & \code{\%r = call i32 @f(i32 \%x)} & i32 \\
无结果调用 & \code{call void @g(i32 \%x)} & 无 \\
无条件跳转 & \code{br label \%next} & 无 \\
条件跳转 & \code{br i1 \%cond, label \%yes, label \%no} & 无 \\
返回整数 & \code{ret i32 \%r} & 无 \\
返回 Unit & \code{ret void} & 无 \\
不可达出口 & \code{unreachable} & 无 \\
phi（拓展） & \code{\%r = phi i32 [ \%a, \%left ],}\newline
  \code{[ \%b, \%right ]} & i32 \\
\end{longtable}
}

\textbf{表外前提也要检查。}移位计数按项目规则屏蔽；除法的除数不能为 0，有符号最小值除以 -1 的情况也不能直接当成普通有效运算。当前 Rx 规范将相应异常算术排除在有效测试输入之外，具体范围仍以规范为准。内存地址必须有效，bool 字节按项目约定为 0 或 1。

\lesson{术语、示例索引与继续阅读}{app:reading}
\subsection{最小术语表}
\begin{tabularx}{\linewidth}{@{}p{43mm}Y@{}}
\toprule
术语 & 本讲义中的含义 \\
\midrule
Module / 模块 & 顶层容器，含类型声明、全局对象和函数。\\
Instruction / 指令 & 一步明确的计算、访问、调用或控制流操作。\\
Operand / 操作数 & 指令使用的值，可能是常量、参数或已有结果。\\
SSA value & 有一个静态定义的局部值，不等于可反复赋值的源变量。\\
Basic block / 基本块 & 从标签开始、以终结指令结束的指令序列。\\
Terminator / 终结指令 & 决定该块的出口，例如 br、ret。\\
CFG / 控制流图 & 用块和跳转边表示可能的执行路径。\\
Predecessor / 前驱 & 有控制流边进入当前块的块。\\
Dominance / 支配 & 到达某个位置的所有入口路径都经过另一位置。\\
GEP & 根据对象布局与索引计算地址的指令。\\
ABI & 调用与布局约定，包括参数、返回、大小和对齐等。\\
Lowering / 降低 & 将较高层语义翻译成较低层表示的过程。\\
Poison & 错误承诺或某些无效操作产生的特殊值，后续使用可能导致未定义行为；不是一个普通的随机整数。\\
\bottomrule
\end{tabularx}

\subsection{配套示例索引}
\begin{tabularx}{\linewidth}{@{}p{50mm}Yp{36mm}@{}}
\toprule
文件（均位于 examples） & 学习重点 & 预期输出 \\
\midrule
\code{01\_print42.ll} & 模块与入口包装 & 42 \\
\code{02\_add.ll} & 参数、call、返回值 & 42 \\
\code{03\_variables.ll} & alloca/load/store & 42 \\
\code{04\_if\_slots.ll} & 分支与合流槽 & 42 \\
\code{05\_short\_circuit.ll} & 短路及副作用 & 0、999、1 \\
\code{06\_loop.ll} & 循环与零次迭代 & 55、0 \\
\code{07\_arrays.ll} & 数组与动态下标 & 42 \\
\code{08\_structs.ll} & 字段布局、bool & 42、1 \\
\code{09\_phi.ll} & SSA 合流拓展 & 42 \\
\code{10\_references.ll} & 指针槽与所指对象 & 42 \\
\code{11\_aggregate\_return.ll} & 隐藏输出地址 & 42 \\
\code{12\_scalar\_rules.ll} & 回绕、比较、右移 & 0、1、0、-4、2147483644 \\
\bottomrule
\end{tabularx}
表中多个数字分别占一行；精确输出以同名 \code{.out} 文件为准，最后一行也有换行。

\subsection{官方资料怎样读}
以下链接已按 LLVM 22.1.0 固定版本。建议“先用小例子建立认识，再按需查手册”，不必从头背完 LangRef。

\begin{enumerate}
 \item \textbf{第一份教程：}
 \href{https://releases.llvm.org/22.1.0/docs/tutorial/MyFirstLanguageFrontend/LangImpl03.html}{Kaleidoscope 第 3 章：Code generation to LLVM IR}。
 重点看 Value、Function、BasicBlock、IRBuilder 之间的关系，以及常量、二元运算和调用的生成。教程语言主要用浮点数，本项目使用整数，不要直接照抄 fadd 等选择。
 \item \textbf{可变变量：}
 \href{https://releases.llvm.org/22.1.0/docs/tutorial/MyFirstLanguageFrontend/LangImpl07.html}{Kaleidoscope 第 7 章：Mutable Variables}。
 重点看为什么 alloca/load/store 适合基础实现，以及 LLVM 后续如何将适合的栈变量提升成 SSA 值。优化章节可后读。
 \item \textbf{控制流：}
 \href{https://releases.llvm.org/22.1.0/docs/tutorial/MyFirstLanguageFrontend/LangImpl05.html}{Kaleidoscope 第 5 章：Control Flow}。
 重点看 if、循环、基本块切换和 phi。理解其 CFG 思路，再适配项目目前的合流槽方案。
 \item \textbf{地址专项：}
 \href{https://releases.llvm.org/22.1.0/docs/GetElementPtr.html}{The Often Misunderstood GEP Instruction}。
 重点看第一个索引、数组和结构体的区别，以及为什么 GEP 不进行内存解引用。
 \item \textbf{新旧资料衔接：}
 \href{https://releases.llvm.org/22.1.0/docs/OpaquePointers.html}{Opaque Pointers}。
 看懂旧例子中的 \code{i32*} 与新例子中 \code{ptr} 的差异，明确类型信息现在写在 load/store/GEP 等操作上。
 \item \textbf{权威语法查询：}
 \href{https://releases.llvm.org/22.1.0/docs/LangRef.html}{LLVM Language Reference Manual}。
 遇到一条指令时，查它的语法、参数要求和语义；尤其核对 align、inbounds、nsw 等附加承诺。先读 Integer Type、Function、Basic Block 及本讲义出现的指令即可。
 \item \textbf{源语言边界：}
 \href{https://github.com/ACMClassCourse-2025/rx-compiler-specification/blob/master/src/backend.md}{Rx 规范：后端目标}，
 \href{https://github.com/ACMClassCourse-2025/rx-compiler-specification/blob/master/src/expressions/operator-expr.md}{Rx 规范：运算符表达式}。
 它们回答“本项目需要什么行为”；本仓库的 \code{docs/ir.md} 与 \code{docs/semantic.md} 则说明如何在现有实现中组织这些信息。
\end{enumerate}

\textbf{暂时可以放后。}调试元数据、异常处理、原子指令、向量指令、内联汇编、自定义 pass、复杂优化属性和完整 SSA 构造，都不是打印、标量、控制流和基础聚合阶段的前置条件。

\textbf{资料范围。}本讲义是基于官方语义和当前项目约定编写的中文教学说明，不是 LLVM 或 Rx 规范的完整替代。LLVM 22 的通用语法与项目选择已分别标注；项目代码继续演进时，接口细节以仓库实际版本为准。

\end{document}
